%=== CHAPTER SIX ===
\chapter{Conclusion and Future Work}
\label{ch:conclusion}
\begin{spacing}{1.5}
\setlength{\parskip}{0.22in}

\section{Summary}

This dissertation has structured the quadrotor payload problem as a coupled estimation-and-control task. The predictive model changes at grasp and release in three ways: mass changes translational acceleration and hover thrust, attachment geometry changes inertia, and horizontal CoM offset creates a trim moment. Treating these effects as one undifferentiated disturbance obscures their different physical channels.

The proposed architecture estimates composite mass as a fourteenth MHE
state over a \SI{2}{s} sliding window without inserting the
empty-airframe or parcel mass into its initialisation, arrival weight or
lower bound. Physical thrust and torque are reconstructed from rotor
speed to avoid self-confirming use of the NMPC's intended input. A
pre-offboard hover establishes an unloaded estimate; after pickup the
parcel prediction is $\hat m_p=(\hat m_t-m_b)^+$. When a
discontinuity is confirmed, pre-event stages are deweighted instead of
forcing one constant mass to compromise between two physical regimes. A
signed thrust-step detector identifies pickup and drop. At drop it
releases obsolete geometry but never resets mass, leaving the MHE to
regress from the loaded estimate to the unloaded solution. Rigid-body
relations provide the payload-induced inertia increment, while horizontal
CoM can be inferred from steady torque balance. Revision \texttt{f2264b3}
adds a bounded online-$\dJ$ tracker so that the published mass prediction
is actually propagated into the NMPC inertia model. NMPC then consumes one
consistent parameter set in both its dynamics and hover-input baseline.

The complete pipeline is implemented around ROS~2, PX4 SITL, Gazebo,
MAVROS and acados. The MHE operates at \SI{10}{Hz}, while the 20-stage
NMPC solves at \SI{20}{Hz} over a \SI{1}{s} preview. A
three-dimensional figure-eight trajectory, direct mass-step plugin and
detachable magnetic gripper provide increasingly coupled validation
conditions. The repository-aligned experiment matrix includes prior-free
cold-start observability, unknown parcel-mass prediction, drop regression,
a residual-only trigger and an L1-adaptive baseline.

\section{Principal Findings}

At revision \texttt{f2264b3}, the defensible algorithmic findings are:
\begin{enumerate}
    \item removing a mass label requires deleting all three information paths: the mass-derived lower bound, the mass arrival weight and the nominal-mass fallback;
    \item rotor thrust supplies a model-independent initial iterate and measurement channel, whereas intended thrust is self-confirming;
    \item parcel mass can be reported as the optimised total-mass increment without becoming a separate estimator state or an input prior;
    \item pickup/drop direction is observable from the sign of a persistent physical-thrust step, although trajectory switches remain a known ambiguity;
    \item geometry release and mass-state reset are different operations: only the former is permitted at drop, so unloaded recovery remains an estimation result;
    \item the NMPC-side inertia estimate can follow the predicted payload mass through a bounded ratchet or a two-way tracker, but the floor, cap and gain-envelope settings must be reported as safety assumptions;
    \item payload adaptation must update the equilibrium-input penalty as well as the state equation;
    \item predecessor revision \texttt{e003d339} reports event-settling and NMPC timing improvements, but those values require repetition under the new prior-free configuration.
\end{enumerate}

\draftnote{Replace the prior-free acceptance criteria with the frozen
cold-start and drop-regression results, including run identifiers,
confidence intervals and failed-run counts. Keep predecessor revision
e003d339 statistics visually separate.}

\section{Limitations}

The present model assumes rigid attachment and known contact geometry.
The reusable airframe model still contains calibrated $m_b$, inertia and
rotor coefficients; \emph{no prior mass} refers to the MHE state and
unknown parcel, not to removal of all vehicle calibration. Moreover, the
current stage-D controller retains a \SI{0.15}{kg} inertia-map floor and
a \SI{0.5}{kg} platform envelope for attach-time rate-gain scaling. These
values do not enter the MHE or reveal the selected parcel's truth, but
they are still prior information about the admissible payload scale.
Accordingly, strict ``no payload-mass prior anywhere in the closed loop''
has not yet been demonstrated. Thrust-only
event detection cannot fundamentally distinguish a payload step from
every possible sustained vertical-force or trajectory transition, so
the higher geometry-release threshold is a safety separation rather than
a proof of unique attribution. The model also neglects flexible,
swinging and sloshing payload modes, uses a static quadratic thrust map
and omits the yaw inertia increment. Finally, SITL cannot reproduce every
latency, vibration and calibration error of a physical airframe.

\section{Future Work}

Immediate future work is to freeze and archive the evidence package for
the no-task-mass cold-start, pickup prediction, online-inertia and drop-regression grid,
regenerate every table from retained logs and report run-level confidence
intervals. The actuator limits must remain fixed independently of the
estimate. The next modelling step is to include the yaw inertia
increment, propagate uncertainty from rotor-speed calibration into the
MHE weights, run zero-floor/envelope-free controller ablations and replace the simplified arrival cost with a
covariance-consistent update.

Hardware transfer should proceed in stages: motor stand calibration, tethered hover, centred payload pickup, eccentric payload pickup and only then dynamic figure-eight flight. A safety supervisor should monitor solve status, state-estimation quality and input margin. Longer-term extensions include explicit flexible-load dynamics, multi-rate MHE, chance-constrained NMPC under wind and joint estimation of slowly varying aerodynamic coefficients.

\section{Closing Remark}

The main lesson is that online estimation is useful only when the estimated quantity enters the controller in a physically and numerically consistent way. By aligning measurements, parameter structure, equilibrium input and constraints, the MHE--NMPC architecture provides a clear route from payload detection to continued constrained flight.

\end{spacing}
\newpage
